% !TEX program = xelatex
% 通信学报风格模板（通用版）
\documentclass[twocolumn,a4paper,10pt]{ctexart}

% 页面设置（符合通信学报要求）
\usepackage[top=2.2cm,bottom=2.2cm,left=1.8cm,right=1.8cm]{geometry}
\setlength{\columnsep}{0.8cm}

% 字体与段落
\usepackage{xeCJK}
\setCJKmainfont{SimSun}[AutoFakeBold=true,AutoFakeSlant=true] % 宋体
\setCJKsansfont{SimHei} % 黑体
\setmainfont{Times New Roman}
\linespread{1.2} % 1.2倍行距

% 数学与图表
\usepackage{amsmath,amssymb}
\usepackage{graphicx}
\usepackage{float}
\usepackage{booktabs}
\usepackage{cite}

% 标题格式
\ctexset{
  section = {
    format = \heiti\zihao{4}\bfseries\raggedright,
    beforeskip = 10pt,
    afterskip = 5pt
  },
  subsection = {
    format = \heiti\zihao{5}\bfseries,
    beforeskip = 5pt,
    afterskip = 3pt
  }
}

% 摘要环境重定义
\renewenvironment{abstract}{
  \begin{center}
  \begin{minipage}{0.9\linewidth}
  \zihao{5}\kaishu
  {\heiti\zihao{5} 摘\quad 要：}
}{
  \end{minipage}
  \end{center}
  \vspace{5pt}
}

% 关键词
\newcommand{\keywords}[1]{
  \noindent{\heiti\zihao{5} 关键词：}\zihao{5}\kaishu #1\par
  \vspace{8pt}
}

% 作者信息
\newcommand{\authorinfo}[2]{
  \begin{center}
    {\zihao{4}\fangsong #1}\\[3pt]
    {\zihao{5}\songti (#2)}
  \end{center}
  \vspace{5pt}
}

% 中图分类号等
\newcommand{\clc}[4]{
  \noindent\zihao{6}\songti
  中图分类号：#1 \quad 
  文献标志码：#2 \quad 
  文章编号：#3 \quad 
  DOI：#4\par
  \vspace{5pt}
}

% 标题和作者信息
\title{技术论文示例：基于深度学习的通信优化}
\author{张三, 李四, 王五}
\date{2026年4月26日}

\begin{document}

\maketitle

\begin{abstract}
本文提出了一种基于深度学习的通信优化方法。通过实验验证，该方法在性能和效率上均优于传统方法。
\end{abstract}

\section{引言}
近年来，深度学习在通信领域的应用受到了广泛关注。本文旨在探讨如何利用深度学习技术优化通信系统。

\section{相关工作}
在通信优化领域，传统方法主要包括基于规则的优化和启发式算法。然而，这些方法在复杂场景下的表现有限。

\section{方法}
本文提出了一种基于深度学习的通信优化框架。该框架包括以下几个模块：
\begin{itemize}
    \item 数据预处理：对原始通信数据进行清洗和特征提取；
    \item 模型训练：使用深度神经网络对数据进行建模；
    \item 性能评估：通过仿真实验验证模型的有效性。
\end{itemize}

\section{实验结果}
实验结果表明，本文方法在吞吐量和延迟方面均优于基线方法。具体结果如表\ref{tab:results}所示。

\begin{table}[ht]
    \centering
    \caption{实验结果对比}
    \label{tab:results}
    \begin{tabular}{|c|c|c|}
        \hline
        方法 & 吞吐量 (Mbps) & 延迟 (ms) \\
        \hline
        基线方法 & 50 & 100 \\
        \hline
        本文方法 & 70 & 80 \\
        \hline
    \end{tabular}
\end{table}

\section{结论}
本文提出了一种基于深度学习的通信优化方法，并通过实验验证了其有效性。未来工作将进一步优化模型结构，并探索更多应用场景。

\bibliographystyle{ieeetr}
\bibliography{references}

\end{document}